\chapter{Discussion and Conclusions}\label{sec:discussion-and-conclusions}

\paragraph{Scope of the discussion.} The conclusions discussed in this chapter draw on two distinct parts of the thesis. Chapter~\ref{sec:experiments} evaluates relay strategies under the canonical memoryless channel model, where transmission is uncoded, the noise is white, and processing is performed on a per-symbol basis. Chapter~\ref{sec:unknown-channels} then deliberately departs from that model by introducing impairments such as intersymbol interference (ISI), nonlinear distortion, and unknown phase. Accordingly, statements concerning channel memory, sequence detection, Viterbi MLSE, pilot-based estimation, or blind equalization should be understood as applying only to the extension study of Chapter~\ref{sec:unknown-channels}, not to the canonical memoryless setup.

The canonical core (Chapter~\ref{sec:experiments}) settles hypotheses H1--H4 and establishes the baseline finding that, on a matched known channel, a minimal learned relay does not improve on classical decode-and-forward. The thesis's principal contribution (Chapter~\ref{sec:unknown-channels}) settles H6 and is where the substantive new claim lies: a systematic delineation of when a learned relay adds value under unknown and mismatched channels. The interpretation below treats the two on equal footing, with the unknown-channel results carrying the main argument.

\section{Interpretation of Results}\label{sec:interpretation-of-results}

The experimental results reveal several consistent patterns across the canonical setup and its two scoped extensions. This section interprets these patterns through the lens of the theoretical framework established in Section \ref{sec:theoretical-foundations-channel-models}  and evaluates each research hypothesis.

\subsection{Low-SNR Advantage of Neural Relays over AF (H1: Confirmed)}\label{sec:low-snr-advantage-of-neural-relays-h1-confirmed}

\textbf{Low SNR (0-4 dB): AI advantage over AF.} At low SNR, AI-based relays consistently outperform AF on the canonical channel, confirming H1. The theoretical explanation lies in the structure of the MMSE-optimal (posterior-mean) relay estimate. After coherent compensation, each axis of the canonical QPSK-over-Rayleigh link reduces to the same real-valued single-sample problem as BPSK over AWGN (Section~\ref{sec:two-hop-relay-model}); the derivation below is carried out in that equivalent real form, applied identically per axis:

\begin{equation}\protect\phantomsection\label{eq:optimal-denoiser-relay}{f^*(y) = \mathbb{E}[x | y] = \tanh\left(\frac{y}{\sigma^2}\right)
}\end{equation}
\addcontentsline{equ}{listequations}{\protect\numberline{\theequation}\quad optimal-denoiser-relay}

At low SNR (e.g., 0 dB, \(\sigma^2 = 1\)), this is a gentle sigmoid: \(f^*(y) = \tanh(y)\). The neural network can approximate this smooth function accurately, producing a \textbf{soft estimate} that preserves information about the confidence of the decision. By contrast:

\begin{itemize}
\item
  \textbf{DF} applies the hard-decision function \(f_{\text{DF}}(y) = \text{sign}(y)\), which discards all soft information. At low SNR, many received samples lie near the decision boundary (\(|y| \approx 0\)), and DF assigns them full confidence (\(\pm 1\)) regardless of the actual uncertainty. This information loss leads to excess errors on the second hop.
\item
  \textbf{AF} applies \(f_{\text{AF}}(y) = Gy\), which is linear and preserves the noisy signal structure but amplifies the noise by the same factor as the signal. The effective SNR degradation is given by \(\text{SNR}_{\text{eff}} = \gamma^2 / (2\gamma + 1) \approx \gamma/2\) at high SNR: a 3 dB penalty.
\end{itemize}

The AI relay implements a \textbf{non-linear soft mapping} that is intermediate between these extremes: it suppresses noise (like DF's regeneration) while preserving soft information near the decision boundary (unlike DF's hard decision). This explains why learned non-linear relay processing can outperform AF and, under selected channel conditions, also outperform DF at low SNR.

Among the AI methods, the low-SNR advantage over AF is realized by the \emph{minimal} feedforward relays, not by the larger sequence models. On the canonical Rayleigh channel (Table~\ref{tbl:table2}, 0~dB) the Hybrid, MLP and VAE relays reach $0.3329$, $0.3349$ and $0.3359$ respectively, essentially matching symbol-wise DF's $0.3337$ and improving substantially on AF's $0.4076$, while the higher-capacity Transformer, Mamba-S6 and Mamba-2 (SSD) relays sit at $0.4014$, $0.4071$ and $0.4032$ -- that is, at essentially AF's level, capturing almost none of the available advantage. The advantage over AF is therefore \emph{not} driven by parameter count: the smallest network already captures it, consistent with the complexity ordering discussed under H3 and with the convergence observed in the equal-parameter-budget comparison (H4, Section~\ref{sec:normalized-parameter-comparison}).
\REV{Fix applied (contradiction with Table~\ref{tbl:table2}): this paragraph previously claimed ``Mamba S6 achieves the lowest BER \ldots ordering correlates with capacity'', which the canonical-channel table contradicts. Second correction: the replacement text still carried stale BER values ($\approx0.247$/$0.245$/$0.317$--$0.325$) predating the SNR-convention correction, and named the cGAN, which is excluded from the reported comparison. All figures now transcribed directly from Table~\ref{tbl:table2}'s 0~dB row; the claim that the VAE was ``worse still'' was also false and has been removed.}

\subsection{Classical Dominance at High SNR (H2: Confirmed)}\label{sec:classical-dominance-at-high-snr-h2-confirmed}

\textbf{Medium-to-high SNR (≥6 dB): Classical dominance.} At medium and high SNR, DF matches or exceeds all AI methods with exactly zero parameters and zero training time. The theoretical explanation is straightforward: at high SNR, \(\sigma^2 \to 0\) and \(f^*(y) \to \text{sign}(y) = f_{\text{DF}}(y)\). The symbol-wise DF relay \emph{exactly implements} the limiting form of the posterior-mean estimate in this regime, while any neural network approximation introduces a small but non-zero reconstruction error due to the finite precision of the learned mapping and the tanh output saturation (which approaches but never reaches \(\pm 1\)).

The effect is visible analytically in the AWGN calibration limit, at 10 dB the two-hop DF BER is \(2Q(\sqrt{10}) \cdot (1 - Q(\sqrt{10})) \approx 0.002\), and empirically on the canonical Rayleigh channel, where DF attains the lowest or tied-lowest BER at every point from 6 dB upward. The AI models cannot beat DF in this regime, for the uncoded per-symbol setting, $\operatorname{sign}(y)$ is the MAP symbol decision and hence optimal among per-symbol relay functions (Remark~\ref{rem:df-terminology}), they can only match it, and the larger sequence models slightly underperform due to the increased variance of their more complex learned mappings.

\subsection{Robustness Across the Fading Ensemble}\label{sec:channel-robustness}

\textbf{Robustness across fades.} The relay ranking on the canonical channel is stable across the full evaluated SNR range, even though each transmitted symbol experiences an independent fading realization. The mechanism is that after coherent compensation (\(\hat{x} = y \cdot h^*/|h|\)), the relay processes a signal with AWGN-like noise at a \emph{random effective SNR}; the network, trained across multiple SNR values, has learned a denoising function that adapts to different noise levels, so the Rayleigh ensemble (a mixture of effective SNRs weighted by the fading distribution) is handled by virtue of multi-SNR training alone, with no per-realization retraining or explicit CSI input.

\section{The ``Less is More'' Principle (H3: Partially Supported)}\label{sec:the-less-is-more-principle-h3-confirmed}

One of the most practically useful findings is that relay BER does not improve with model size over the range evaluated: the Minimal MLP architecture (169 parameters) matches the performance of models with 100--140$\times$ more parameters (3K--24K), and no larger evaluated model improved on it under the training protocol used. The stronger reading of H3 --- that the curve turns upward because excess capacity causes \emph{overfitting} --- is not established by the evidence collected here. Each model was trained once, with no seed replication, no train/test-gap measurement and no held-out validation curve, so a size effect cannot be separated from run-to-run variance or from an optimization difficulty that more capacity happens to introduce. The claim carried forward is therefore the weaker, directly measured one: larger evaluated models did not improve BER. What follows in this section is offered as a candidate explanation for that observation, not as a demonstrated mechanism.
\REV{``inverted-U relationship between model complexity and relay performance'' \(\to\) ``U-shaped \ldots in BER (equivalently, an inverted-U in accuracy)'', consistent with H3 as retitled in Chapter~\ref{sec:research-objectives}.}

This result has important theoretical and practical implications:

\subsection{Information-Theoretic Perspective}\label{sec:information-theoretic-perspective}

The canonical memoryless relay-denoising task maps a window  of \(2w+1\) real-valued noisy observations to a single real-valued clean estimate, applied per axis to the canonical QPSK symbol (equivalently, on the real BPSK/AWGN calibration channel). Each real axis carries exactly 1 bit of information, and the Bayes-optimal estimator \(f^*(\mathbf{y}) = \tanh(\mathbf{1}^T \mathbf{y} / \sigma^2)\) (for i.i.d. noise across the window) is a one-dimensional function of the sufficient statistic \(\sum_i y_i\). The \textbf{effective dimensionality} of this mapping is therefore 1: far below the \(2w+1 = 5\) or \(11\) input dimensions.

As a heuristic, the minimum description length (MDL) principle suggests matching model complexity to the effective dimensionality of the target mapping. Since the target here is a one-dimensional monotone function of the sufficient statistic, even the 169-parameter network has capacity far in excess of what the mapping requires; this is offered as an intuition for why the smallest model suffices, not as a formal MDL bound.

\subsection{Bias-Variance Analysis}\label{sec:bias-variance-analysis}

The bias-variance decomposition \cite{GoodfellowBengioCourville2016DeepLearning} provides a quantitative framework:

\begin{equation}\protect\phantomsection\label{eq:bias-variance-relay}{\text{MSE}(\hat{x}) = \underbrace{(\mathbb{E}[\hat{x}] - f^*(y))^2}_{\text{Bias}^2} + \underbrace{\mathbb{E}[(\hat{x} - \mathbb{E}[\hat{x}])^2]}_{\text{Variance}} + \underbrace{\sigma^2_{\text{irred}}}_{\text{Noise floor}}
}\end{equation}
\addcontentsline{equ}{listequations}{\protect\numberline{\theequation}\quad bias-variance-relay}

\begin{itemize}
\tightlist
\item
  \textbf{169 parameters:} Low bias (sufficient capacity for the simple \(f^*\)) and low variance (limited capacity prevents memorization). The model operates at the optimal bias-variance trade-off.
\item
  \textbf{3,000 parameters:} Similar bias (the target function hasn't changed) and slightly higher variance. Performance is nearly identical to 169 params.
\item
  \textbf{24,001 parameters (Mamba-S6):} Despite having 140$\times$ more parameters than the MLP, it does not improve on it; on the canonical channel it is worse at low SNR (Table~\ref{tbl:table2}). The marginal utility of the additional parameters is at best vanishingly small.
\end{itemize}

This decomposition is a framework for interpreting the observation, not a measurement of it: bias and variance were not estimated separately, which would require training each configuration repeatedly across seeds and comparing training against held-out error. An earlier version of this section additionally attributed ``clear overfitting'' to an 11{,}201-parameter Maximum MLP; that configuration is not part of the reported experiments and no such measurement exists, so the claim has been withdrawn rather than restated.

\subsection{Practical Implications}\label{sec:practical-implications}

\begin{enumerate}
\def\labelenumi{\arabic{enumi}.}
\item
  \textbf{Occam's Razor for relay AI.} The canonical relay denoising task has inherently low complexity: the per-axis mapping from noisy to clean symbols is fundamentally simple, and the neural network needs only enough capacity to learn a non-linear threshold function with some contextual smoothing. Adding parameters beyond this minimal requirement bought no measurable accuracy anywhere in this thesis.
\item
  \textbf{Deployment feasibility.} A 169-parameter model requires approximately 0.7 KB of memory (169 float32 values) and trains in seconds on a CPU. This makes AI-based relay processing viable even on severely resource-constrained embedded relay nodes, such as IoT devices or sensor network relays. The model can be stored in on-chip SRAM and executed without any external memory access.
\item
  \textbf{Generalization.} The 169-parameter MLP handles the full Rayleigh fading ensemble, every effective SNR the fades produce, with a single set of weights and no per-condition retraining, an instance of generalization that would be harder to achieve with a larger, more specialized model.
\end{enumerate}

\section{Practical Deployment Recommendations}\label{sec:practical-deployment-recommendations}

Based on the comprehensive evaluation, we propose the following deployment strategy:

\begin{longtable}[]{@{}
  >{\raggedright\arraybackslash}p{(\linewidth - 4\tabcolsep) * \real{0.3333}}
  >{\raggedright\arraybackslash}p{(\linewidth - 4\tabcolsep) * \real{0.3333}}
  >{\raggedright\arraybackslash}p{(\linewidth - 4\tabcolsep) * \real{0.3333}}@{}}
\caption[Practical deployment recommendations]{Practical deployment recommendations: recommended relay strategy by operating regime.}\label{tbl:table32}\tabularnewline
\toprule\noalign{}
\begin{minipage}[b]{\linewidth}\raggedright
Operating Regime
\end{minipage} & \begin{minipage}[b]{\linewidth}\raggedright
Recommended Relay
\end{minipage} & \begin{minipage}[b]{\linewidth}\raggedright
Rationale
\end{minipage} \\
\midrule\noalign{}
\endfirsthead
\toprule\noalign{}
\begin{minipage}[b]{\linewidth}\raggedright
Operating Regime
\end{minipage} & \begin{minipage}[b]{\linewidth}\raggedright
Recommended Relay
\end{minipage} & \begin{minipage}[b]{\linewidth}\raggedright
Rationale
\end{minipage} \\
\midrule\noalign{}
\endhead
\bottomrule\noalign{}
\endlastfoot
\textbf{Low SNR (0-4 dB)} & MLP Minimal / Hybrid / channel-specific selection & Neural relays often help; exact best choice is channel-dependent, with DF remaining strong on some fading scenarios \\
\textbf{Medium SNR (4-8 dB)} & Hybrid (MLP + DF) & Automatic switching at empirically chosen crossover\\
\textbf{High SNR (\textgreater8 dB)} & DF & Zero parameters; lowest BER among the evaluated per-symbol relays \\
\textbf{Resource-constrained} & MLP Minimal (169 params) & 0.7 KB memory, \textless5 s training \\
\textbf{Best overall (BPSK/QPSK)} & Hybrid & Combines AI advantage with DF reliability \\
\end{longtable}
\REV{``optimal performance'' \(\to\) ``lowest BER among the evaluated per-symbol relays'' in the High-SNR row, to avoid an unbounded optimality claim (consistent with the H2 wording used elsewhere).}

The Hybrid relay is recommended as the default deployment choice because it automatically selects MLP processing at low SNR and DF at high SNR, achieving near-optimal performance across the entire SNR range with minimal complexity.

\subsection{Latency and Adaptivity Considerations}\label{sec:latency-and-adaptivity-considerations}

Two practical questions bear on whether a learned relay can be deployed with minimal impact: the latency it adds relative to AF/DF, and its ability to adapt to varying channel conditions without retraining.

\textbf{Latency.} Two distinct latency components must be separated. The first is \emph{buffering} latency, incurred by the sliding window $y_R[i-w:i+w]$. In the half-duplex protocol studied here this component is absorbed by the protocol itself: the relay buffers the entire listening-phase block before the transmission phase begins (Remark~\ref{rem:window-causality}), so the window's $w$-symbol look-ahead adds no delay beyond what store-and-forward operation already imposes. Only in a streaming or ultra-low-latency protocol, outside the scope of this thesis, would the symmetric window cost $w$ symbols of look-ahead that the memoryless AF/DF baselines ($w=0$) do not pay. The second component is \emph{compute} latency: inference cost per relayed symbol. Here the U-shaped complexity result (H3) is decisive: since the 169-parameter MLP matches the performance of models $100-140\times$ larger, the deployed network reduces to two small matrix--vector products per symbol, a per-symbol cost that is orders of magnitude below a single channel use's duration on any modern embedded processor, and negligible in absolute terms even if it cannot literally match the single comparison of DF's $\operatorname{sign}(y)$. The compute-latency argument therefore holds \emph{because of}, not despite, the less-is-more finding; it would not hold for the 24K-parameter sequence models, whose recurrent or attention-based structure imposes materially higher per-symbol cost for no BER benefit at the relay window length.
\REV{``the inverted-U result (H3)'' \(\to\) ``the U-shaped complexity result (H3)''.}

The coded study of Section~\ref{sec:coded-latency-throughput} adds a third component that the canonical comparison never had to consider, and it is the largest of the three: \emph{buffering}. A block relay cannot emit until it has decoded a whole frame ($202$ symbols there), and the soft BCJR variant cannot even in principle, its backward recursion running from the end of the frame. A windowed learned relay needs only its $w$-symbol look-ahead, ten symbols, and that figure is constant in frame length where the block relays' grows with it. The measured arithmetic points the same way ($0.30\,\mu$s per symbol against Viterbi's $15.07$), so on the coded link the learned relay's advantage in delay and cost is structural rather than marginal, even where its BER advantage is small or absent.

\textbf{Adaptivity.} A single network, trained across the range of simulated SNRs, handles the entire canonical fading ensemble without retraining or explicit CSI input (Section~\ref{sec:channel-robustness}): after coherent phase compensation, each fading realization presents the relay with an AWGN-like observation at a random effective SNR, and multi-SNR training already covers this mixture. Two honest qualifications apply. First, this adaptivity is an advantage over AF, whose fixed gain is tuned to an assumed SNR, but not over symbol-wise DF, whose $\operatorname{sign}(y)$ slicer is equally condition-agnostic; adaptivity therefore does not alter the H2 conclusion. Second, adaptation has been demonstrated only within the canonical Rayleigh ensemble; generalization to other channel families (Rician, bursty interference, phase noise, hardware non-linearities) remains untested and is identified as future work.

\section{Limitations}\label{sec:limitations}

Several limitations of this study should be acknowledged, as they define the boundary conditions under which the conclusions hold:

\begin{enumerate}
\def\labelenumi{\arabic{enumi}.}
\item
  \textbf{Modulation scope.} All experiments use QPSK, the canonical constellation. Higher-order constellations are not investigated: 16-QAM and denser grids break the per-axis relay formulation and require a joint $N$-class 2D relay whose output dimension scales with $M$, which is left to future work.
\item
  \textbf{Perfect CSI.} We assume perfect channel state information at each receiver. In practice, channel estimation errors would affect the coherent compensation step and hence relay performance. Imperfect CSI introduces a model mismatch between the assumed and actual channel, which could differentially impact the relay strategies: AI relays might be more robust (having learned from noisy data) or less robust (having not been exposed to estimation errors during training).
\item
  \textbf{Static offline training.} AI relays are trained once on synthetic data and deployed without adaptation. While multi-SNR training covers the canonical fading ensemble (Section \ref{sec:channel-robustness}), and the unknown-channel study (H6) shows a fixed network can learn a \emph{stationary} unknown impairment, this protocol does not exploit online information; adaptive or continual training would be required if the deployment channel drifts over time.
\item
  \textbf{Single relay.} The framework considers a single relay node. Multi-relay cooperation introduces relay selection, power allocation, and cooperative diversity dimensions not captured in the two-hop model. With multiple relays, the system-level optimization (which relay to use, how to combine multiple relay observations) becomes as important as the per-relay processing function studied here.
\item
  \textbf{Two-hop only.} Extension to multi-hop networks with 3+ relays introduces noise accumulation (for AF) and error propagation (for DF) that grow with the number of hops. AI relays might show greater advantage in multi-hop settings where noise accumulation is more severe.
\item
  \textbf{Fixed window size.} The relay window size (\(w = 2\) or \(w = 5\)) is fixed and relatively small. Larger windows could provide more context for denoising, particularly in channels with temporal correlation (e.g., block fading). A systematic study of window-size optimization is deferred.
\item
  \textbf{No coding, in the core comparison.} The core comparison operates without error-correcting codes, so the DF baseline there is its symbol-level variant (Remark~\ref{rem:df-terminology}). Sections~\ref{sec:coded-block-df}--\ref{sec:coded-latency-capacity} measure a coded, information-theoretic block-DF baseline, its soft-decision counterparts, a rate-adaptive envelope over a modulation-and-coding grid, and that envelope re-derived under a round-trip latency budget, directly rather than leaving any of this as an assertion, but that study is confined to one code family, one frame length, and one channel model, and the relay and destination each decide once with no iterative exchange between them.
\item
  \textbf{Single training run per model; unequal training budgets across architectures.} Each learned relay is trained once, from one random initialization, rather than across multiple seeds with mean/standard-deviation/confidence-interval reporting of training variance; the confidence intervals reported throughout this thesis are Monte Carlo intervals over the \emph{evaluation} channel and bit realizations for a fixed trained model, not over the training procedure itself. Separately, the architectures do not all receive identical training budgets (e.g.\ the cGAN trains for 200 epochs against 100 for the other relays, Section~\ref{sec:relay-strategies}), a confound the normalized-parameter study (Section~\ref{sec:normalized-parameter-comparison}) does not remove, since it equalizes parameter count and input window but not epoch count. Consequently, the U-shaped complexity finding (H3) and the claim that architecture is a secondary factor relative to capacity (Section~\ref{sec:parameter-normalization-complexity-trade-off}) should be read as observations from the single trained instance of each architecture studied here, not as claims that have been shown robust to training-seed variance or fully isolated from training-budget differences; a multi-seed re-run with equalized budgets is the direct way to close this gap and is not performed in this thesis.
\end{enumerate}

\section{Future Work}\label{sec:future-work}

Several directions warrant further investigation:

\begin{enumerate}
\def\labelenumi{\arabic{enumi}.}
\item
  \textbf{Time-selective fading.} The present study addresses channels that are static within a block. Extending the learned relay to time-selective fading with continuously varying unknown phase (e.g.\ Jakes/Clarke Doppler) is a natural next step, likely requiring explicitly phase-aware architectures (complex-valued state, learned phase unwrapping, or a training objective aligned with noncoherent detection rather than symbol MSE) evaluated against prediction-based noncoherent receivers.
\item
  \textbf{Denser constellations (64-QAM, 256-QAM).} The extension shows joint $N$-class 2D classification eliminates the 16-QAM floor; scaling the formulation to 64 or 256 classes may require deeper networks or hierarchical (coarse-to-fine) classification, and whether the U-shaped complexity curve (H3) shifts rightward with constellation density is a key open question.
\item
  \textbf{Imperfect CSI.} Introduce channel estimation errors to assess robustness of AI relay processing under realistic conditions.
\item
  \textbf{A BER-optimal (BCJR/APP) benchmark for the unknown-ISI channel.} Section~\ref{sec:unknown-channel-experiment} benchmarks against genie-CSI Viterbi MLSE, the optimal \emph{sequence} detector. Sequence-ML and bit-BER-optimal detection are distinct criteria on any channel with memory, at both modulation orders studied here, so the benchmark is not BER-optimal in either case; for Gray-coded QPSK there is the additional effect that a wrong trellis path can differ from the truth by one or two bits with unequal probability. A BCJR/APP detector run over the same 3-tap trellis, which directly minimizes per-bit error probability rather than sequence error probability, would test whether the classical benchmark's margin over the learned relay changes once it is BER-optimal rather than sequence-optimal, and is the measurement needed to settle the mechanism conjectured in Section~\ref{sec:qpsk-unknown-channel}; this was not run here.
\item
  \textbf{Online learning.} Develop relay strategies that adapt their parameters during operation, tracking time-varying channel conditions.
\item
  \textbf{Multi-relay networks.} Extend to cooperative relay selection and multi-hop routing with AI-optimized relays at each node.
\item
  \textbf{Other channels and MIMO.} Evaluate the canonical conclusions on additional channel families (Rician line-of-sight fading, and AWGN treated as a full operating point rather than the bounded calibration-plus-companion role it has here) and on a multi-antenna second hop, where destination equalization interacts with relay processing and the additivity of the two gains becomes a testable question.
\item
  \textbf{End-to-end learning.} Train the entire communication chain (modulation, relay, demodulation) jointly using autoencoder-based approaches, and compare against the modular classical-plus-learned-relay cascade studied here.
\item
  \textbf{Coded block-DF baseline (measured; open items remain).} Remark~\ref{rem:df-terminology} scopes the core comparison's DF baseline to the uncoded, symbol-level variant. Section~\ref{sec:coded-block-df} measures the coded, information-theoretic block-DF baseline this item once only proposed, together with the soft-information relays (BCJR at the relay, posterior-mean read-out of the learned relay) and an adaptive-rate envelope over a modulation-and-coding grid. The substantive finding is mechanistic: a relay that decodes the code has spent its redundancy before transmitting, so a wrong decode arrives as a valid-but-wrong codeword the destination cannot repair, whereas a relay that only denoises leaves the redundancy intact. Rate selection then dominates the relay-decoding question by an order of magnitude, unless a latency deadline is imposed, at which point block-DF's double buffering reopens the gap. What remains open: one code family, one frame length, one channel model; the soft block-DF relay was given the nominal noise variance rather than per-symbol CSI, which caused its 20~dB mis-calibration; iterative (turbo-style) relay--destination exchange was not evaluated; and the latency analysis uses structural buffer counts rather than a $10^{-5}$-class reliability target, so it does not establish ultra-reliable low-latency performance.
\item
  \textbf{Longer relay windows with state-space architectures.} Future work should explore whether longer relay windows (e.g., 128--512 symbols) combined with a chunk-parallel state-space architecture can improve both BER performance and processing speed simultaneously; the training-time crossover between sequential and chunk-parallel formulations, removed from this thesis's scope, is the natural companion measurement.
\end{enumerate}

\section{Conclusions}\label{sec:conclusions}

This thesis presents a comprehensive comparative study of nine relay strategies: two classical (AF, DF) and seven AI-based (MLP, Hybrid, VAE, cGAN, Transformer, Mamba-S6, Mamba-2 SSD). All evaluated on one canonical setup (SISO on both hops, i.i.d.\ Rayleigh fast fading, complex baseband, QPSK, uncoded BER; BPSK appears only on the separate AWGN calibration channel). The study addresses three identified research gaps (Section \ref{sec:research-gap-and-motivation}) through controlled experiments with statistical rigor (100,000 bits per SNR point, 10 independent trials, Wilcoxon significance testing). The following table summarizes the hypothesis outcomes:

\begin{longtable}[]{@{}
  >{\raggedright\arraybackslash}p{(\linewidth - 4\tabcolsep) * \real{0.3333}}
  >{\raggedright\arraybackslash}p{(\linewidth - 4\tabcolsep) * \real{0.3333}}
  >{\raggedright\arraybackslash}p{(\linewidth - 4\tabcolsep) * \real{0.3333}}@{}}
\caption{Research hypotheses, their statements, and experimental outcomes.}\label{tbl:table33}\tabularnewline
\toprule\noalign{}
\begin{minipage}[b]{\linewidth}\raggedright
Hypothesis
\end{minipage} & \begin{minipage}[b]{\linewidth}\raggedright
Statement
\end{minipage} & \begin{minipage}[b]{\linewidth}\raggedright
Result
\end{minipage} \\
\midrule\noalign{}
\endfirsthead
\toprule\noalign{}
\begin{minipage}[b]{\linewidth}\raggedright
Hypothesis
\end{minipage} & \begin{minipage}[b]{\linewidth}\raggedright
Statement
\end{minipage} & \begin{minipage}[b]{\linewidth}\raggedright
Result
\end{minipage} \\
\midrule\noalign{}
\endhead
\bottomrule\noalign{}
\endlastfoot
H1 & Some AI relays outperform AF at low SNR & \textbf{Confirmed}: the best neural relays beat AF at 0--4 dB on the canonical channel \\
H2 & DF lowest-BER among per-symbol relays at high SNR & \textbf{Confirmed}: DF matches or beats all AI methods at \(\geq 6\) dB with 0 parameters \\
H3 & U-shaped complexity curve (BER) & \textbf{Partially supported}: 169 params matches 24K, and larger evaluated models did not improve BER under the training protocol used; the overfitting mechanism is not established, as each model was trained once and no training-variance or train/test-gap analysis was performed \\
H4 & Architecture convergence at equal parameter budget & \textbf{Confirmed at high SNR, rejected at low}: at 3K params the six architectures span 4.1\% BER at 20~dB but 17.4\% at 0~dB, splitting into feedforward and sequence groups. Equalizing the budget required changing width, depth, state dimension and input window jointly (Table~\ref{tbl:table28}), so this equalizes capacity, not the other architectural degrees of freedom \\
H6 & Memoryless classical relays fail on unknown channels; a minimal MLP mitigates by learning, near the model-aware optimum & \textbf{Confirmed}: on unknown ISI, AF/DF hit the analytic 0.25 floor (DF non-monotonic in SNR) while the 169-parameter MLP reaches BER $<5\times10^{-5}$, within $\approx$1--1.5 dB of genie-CSI Viterbi MLSE and indistinguishable from it under the Rayleigh second hop; the same architecture mitigates a composite ISI$\times$PA$\times$phase cascade with no impairment model, tracking $\approx$2 dB behind pilot-aided Viterbi and matching blind CMA when no posterior exists; on the canonical control it only matches DF (H2), the advantage is bounded to structural model-class mismatch representable at minimal capacity \\
\end{longtable}
\REV{H2/H3 statements in the outcome table aligned with the rest of the thesis: ``DF optimal at high SNR'' \(\to\) ``DF lowest-BER among per-symbol relays at high SNR''; ``Inverted-U complexity curve'' \(\to\) ``U-shaped complexity curve (BER)''. \texttt{\$\&lt;5\ldots\$} restored to \texttt{\$<5\ldots\$}; all \texttt{\&amp;} restored to \texttt{\&}.}

The main conclusions, in order of significance, are:

\begin{enumerate}
\def\labelenumi{\arabic{enumi}.}
\item
  \textbf{AI relays beat AF at low SNR (0-4 dB) but not DF.} On the canonical channel, the best neural relays clearly outperform AF in the low-SNR regime while tracking symbol-wise DF within statistical uncertainty; DF remains lowest or tied-lowest at every SNR point. The theoretical explanation is that neural networks approximate the MMSE-optimal soft-threshold estimate \(f^*(y) = \tanh(y/\sigma^2)\), whose end-to-end benefit over the hard slicer is limited in this uncoded setting.
\item
  \textbf{DF attains the lowest per-symbol BER at medium-to-high SNR (≥6 dB) with zero parameters.} The classical decode-and-forward relay requires no training and no parameters yet achieves the best performance among the evaluated per-symbol relays when channel quality is sufficient for reliable first-hop demodulation. At high SNR, the posterior-mean estimate converges to \(\text{sign}(y)\), which is exactly the symbol-wise DF operation.
\REV{``DF is optimal at medium-to-high SNR'' \(\to\) ``DF attains the lowest per-symbol BER \ldots'', consistent with the bounded H2 claim (optimal only among per-symbol relay functions, not information-theoretically).}
\item
  \textbf{No AI architecture separates itself at original parameter counts.} On the canonical channel the best neural relays (MLP, Hybrid, cGAN) perform within statistical uncertainty of one another, and classical DF matches or exceeds all of them across the SNR range. The selective state space model's \(O(n)\) complexity and adaptive filtering mechanism make it an efficient sequence-model choice, but no architectural inductive bias yields a decisive BER advantage on this task.
\item
  \textbf{Architecture matters less than parameter count at equal scale --- but only at high SNR.} When all models are normalized to approximately 3{,}000 parameters (Table~\ref{tbl:table8}), the spread across the six architectures falls monotonically with SNR: $17.4\%$ at 0~dB, $9.3\%$ at 12~dB, $4.1\%$ at 20~dB. In the high-SNR regime the original claim holds and architecture is close to irrelevant. At low SNR it does not: the six models separate cleanly into a feedforward group (MLP, Hybrid, VAE at $0.336$--$0.342$) and a sequence group (Transformer, Mamba-S6, Mamba-2 at $0.400$--$0.407$), a gap far larger than the confidence intervals, with the feedforward models close to DF and the sequence models close to AF.

  Two points are worth separating here. The split is by \emph{family}, not by capacity, since all six carry the same parameter budget --- so at low SNR architecture does matter, and the sequence models' inductive bias is actively unhelpful on a memoryless channel where there is no temporal structure to exploit. And the VAE is no longer the exception it was: at 3K parameters it sits inside the feedforward group at every SNR ($0.0101$ against MLP-3K's $0.0099$ at 20~dB). The earlier reading, that the VAE was a consistent underperformer, does not survive re-measurement. The cause is not the corrected noise convention, which a controlled A/B rules out: training the same relay on data generated under the old convention and under the corrected one produces curves that agree to within Monte Carlo noise at every SNR. The earlier figures came from an implementation that predates the current relay module, so the discrepancy is one of code rather than of channel calibration. See the discussion of the generative paradigm in Section~\ref{sec:generative-vs.-discriminative-paradigms-for-relay-processing}.
\REV{Fix applied (units drift): ``$\sim$1~dB'' changed to ``$\sim$1\% BER'' to match the same claim in Chapter~\ref{sec:experiments} and the H4 row of the outcome table, which both report a relative-BER gap.}
\item
  \textbf{A 169-parameter network is sufficient for relay denoising.} The Minimal MLP architecture achieves performance comparable to models 100-140$\times$ larger, demonstrating that the relay denoising task has inherently low complexity. The bias-variance analysis explains this: the target function is simple (a soft threshold), so additional parameters increase variance without reducing bias.
\item
  \textbf{The Hybrid relay is the recommended practical choice, given a usable estimate of the operating SNR.} By combining AI processing at low SNR with classical DF at high SNR, the Hybrid relay achieves near-optimal performance across the entire SNR range with only 169 parameters. It selects between the two regimes from an SNR estimated out of the received signal power, so no manual per-regime configuration is needed. The recommendation is conditional in a way the experiments here do not discharge: the switching threshold is fixed at the crossover measured on the training sweep, and neither the accuracy of the run-time SNR estimate, the sensitivity of the BER to a misplaced threshold, nor any online adaptation of it was evaluated. All reported Hybrid results are therefore obtained under an accurate operating-SNR estimate, and the margin over its two constituents is small enough at the crossover that a mis-estimate there would plausibly erase it.
\end{enumerate}



\REV{Fix applied (orphaned conclusions): three conclusions were removed here, (8)~``Mamba S6 is the strongest architecture across the higher-order modulation experiments'', (9)~``No neural relay outperforms DF across the full 48-variant combinatorial search'', and (10)~``Input LayerNorm is universally beneficial'' (citing a nonexistent ``Section 4.13''). None of the supporting experiments exist in this restructured version, 16-PSK, the CSI/LayerNorm ablation, and the 48-variant search were all cut (Appendix~E). The in-scope part of the DF-optimality claim is covered by Conclusion~2 and the modulation-extension finding below.}


It is important to note that sequence-detection methods such as Viterbi MLSE appear only in the unknown-channel study of Chapter~\ref{sec:unknown-channels}, where channel memory is introduced explicitly through a 3-tap FIR/ISI channel. The canonical system model of Chapters~\ref{sec:introduction}--\ref{sec:experiments} remains memoryless and therefore requires only per-symbol processing.

The unknown-channel contribution (Chapter~\ref{sec:unknown-channels}) settles the complementary question of \emph{when the learned relay is not merely competitive but necessary}:

\begin{itemize}
\tightlist
\item
  \textbf{On channels violating the deployed relay's model class, learning restores reliability at near-optimal cost (H6: confirmed).} The flat-channel control isolates the cause: against unknown phase, gain and per-branch asymmetry, all memoryless, the classical relay does not fail and the MLP only matches it (gap $\leq 0.0036$), so parameter uncertainty alone is not what learning mitigates. Failure appears only with unmodeled \emph{memory}: on unknown 3-tap ISI, AF and symbol-wise DF sit at the analytic $0.25$ floor and DF actually worsens as SNR rises, while the same 169-parameter MLP learns a joint equalizer--detector and falls below $5\times10^{-5}$ at 16~dB. The Viterbi benchmark bounds the claim: given the channel family and enough pilots, classical estimate-then-MLSE stays $1$--$1.5$~dB ahead, so the learned relay's value is \emph{family-agnosticism} rather than optimality --- near-MLSE reliability with no identification stage, on a composite ISI/nonlinearity/phase cascade it was never told the composition of, and, with no pilots or channel prior at all, matching blind CMA without its instability or per-block adaptation cost. Sweeping the pilot budget makes the boundary quantitative: pilot-aided Viterbi wins down to roughly ten pilots and then collapses at the identifiability threshold, while the learned relay is invariant to it. The crossover is therefore governed by identifiability and, through pilot overhead, by block length. Section~\ref{sec:unknown-channels} gives the full evidence; the complexity accounting there shows the cost side scales linearly in window and output dimension against the trellis's $M^{L}$.
\end{itemize}

These findings demonstrate that neural network-based relay processing is a viable complement to classical approaches in the low-SNR regime, matching, though not surpassing, the per-symbol-optimal DF baseline everywhere else at negligible computational cost. The overarching insight is that \textbf{model complexity should be matched to task complexity}: for the canonical relay denoising task, minimal architectures suffice, and the choice between neural network paradigms, i.e. feedforward or sequential, matters less than proper model sizing and regularization. The practical recommendation is to deploy a Hybrid relay with a 169-parameter MLP sub-network, which combines the neural low-SNR advantage over AF with DF's zero-parameter reliability at high SNR; for 16-QAM deployments, the relay should be formulated as a joint $N$-class 2D classifier rather than a per-axis I/Q-split denoiser; and on links with unmodeled impairments (impairments outside the canonical memoryless model, including ISI, bias, and hardware nonlinearity), a minimal learned relay delivers near-MLSE reliability without channel identification, where the memoryless classical relays fail outright (H6).

\begin{center}\rule{0.5\linewidth}{0.5pt}\end{center}
